%% ============================================================
%%  2-Page Research Proposal — Master's Seminar
%%  "Does It Pay to Play Nice?"
%%  Applied Sports Research — League of Legends
%%  Author: Gian Senpinar | February 2026
%% ============================================================

\documentclass[11pt, a4paper]{article}

% ---------- Packages ----------
\usepackage[margin=2cm]{geometry}
\usepackage[hidelinks]{hyperref}
\usepackage{booktabs}
\usepackage{amsmath}
\usepackage{microtype}
\usepackage{enumitem}
\usepackage{titlesec}
\usepackage[style=apa, backend=biber, sortcites=true, doi=false, url=false]{biblatex}
\addbibresource{references.bib}

% ---------- Typography — compact ----------
\setlength{\parskip}{4pt}
\setlength{\parindent}{0pt}
\pagestyle{empty}
\titleformat{\section}{\normalsize\bfseries}{}{0em}{}[\vspace{-2pt}]
\titlespacing*{\section}{0pt}{8pt}{2pt}
\setlist{nosep, leftmargin=1.2em}

% ============================================================
\begin{document}
% ============================================================

% ---------- Inline title block ----------
\begin{center}
{\large\bfseries Does It Pay to Play Nice?}\\[2pt]
{\normalsize The Effect of Toxic Behavior and Leadership-Style Communication
on Win Rates in \textit{League of Legends}}\\[4pt]
{\small Gian Senpinar\quad$\cdot$\quad Research Proposal — Master's Seminar,
Applied Sports Research\quad$\cdot$\quad February 2026}
\end{center}

\vspace{4pt}
\hrule
\vspace{6pt}

% ============================================================
\section{Introduction and Research Gap}
% ============================================================

\textit{League of Legends} (LoL) is a 5-v-5 multiplayer online battle arena with over
150~million registered accounts \parencite{riotgames2016}, yet its in-game
communication environment is consistently rated among the most toxic in online
gaming. Research finds disruptive behavior in roughly 70\% of competitive matches
\parencite{martinezdeavila2023}, and an analysis of 18~million matches links deviant
behavior to player churn \parencite{kwak2015}. \textcite{achterbosch2021} provide the
most direct evidence to date: across 291~matches, players without in-game conflicts
won significantly more often.

Despite these findings, two gaps persist. First, no study has modeled the effect of
toxicity on win probability at scale using behavioral data with appropriate controls
for individual skill. Second, the role of \textit{leadership-style communication} —
structured shot-calling, coordinating pings, proactive vision control — has received
no empirical attention, even though decades of organizational research show that
directive and transformational leadership behaviors improve team outcomes
\parencite{bass1985, dedreu2003}. Moreover, \textcite{sarcasm2023} find that positive
communication is dismissed in toxic environments, suggesting a moderating
relationship between the two constructs.

This study addresses the resulting research problem: \textit{To what extent do toxic
behavior and leadership-style communication — observable through behavioral
proxies in public match data — affect a team's probability of winning a ranked LoL
game, and does one moderate the other?}

% ============================================================
\section{Theoretical Anchors}
% ============================================================

The study draws on three frameworks.
\textbf{Social Identity Theory} \parencite{tajfel1979}: each five-player team forms a
temporary in-group; toxicity (blame, contempt) fragments group identity and the
psychological safety required for coordinated play \parencite{edmondson1999}.
\textbf{Transformational Leadership Theory} \parencite{bass1985}: proactive pings and
objective coordination map onto directive/inspirational leadership behaviors.
\textbf{Team Conflict Theory} \parencite{dedreu2003}: relationship conflict (the
dominant form of LoL toxicity) consistently impairs team performance.
Jointly, the theories predict that toxicity harms and leadership communication helps
win probability, with the leadership effect attenuated at high toxicity levels.

% ============================================================
\section{Research Questions and Hypotheses}
% ============================================================

\textbf{RQ1 — Toxicity.} Does higher team-level toxic behavior reduce win
probability?
\begin{itemize}
  \item \textbf{H1:} Teams with a higher Toxicity Score have a significantly lower
  probability of winning ($\beta_1 < 0$, $p<.05$).
\end{itemize}

\textbf{RQ2 — Leadership.} Does leadership-style communication independently
increase win probability?
\begin{itemize}
  \item \textbf{H2:} Teams with a higher Leadership Score win more often
  ($\beta_2 > 0$, $p<.05$), controlling for individual skill.
\end{itemize}

\textbf{RQ3 — Moderation.} Is the leadership effect contingent on toxicity level?
\begin{itemize}
  \item \textbf{H3:} The Toxicity$\times$Leadership interaction is significantly
  negative ($\beta_3 < 0$, $p<.05$): leadership communication is less effective
  when toxicity is high.
\end{itemize}

% ============================================================
\section{Data and Methodology}
% ============================================================

\textbf{Data source.}
All data are retrieved from the \textbf{Riot Games Match-V5 API}
(\url{https://developer.riotgames.com}), which provides per-player statistics
(kills, deaths, assists, gold, CS, vision score), per-player \textit{ping counts by
type} (onMyWay, command, assistMe, danger, push, getBack, enemyMissing,
visionCleared), and match timeline events (objective kills, surrender votes).
In-game chat logs are \textit{not} exposed by the public API; accordingly, all
constructs are operationalized through behavioral proxies.

\textbf{Sample.}
2,000 ranked solo/duo matches from the EUW server, stratified across Gold and
Platinum/Emerald tiers. Matches under 15~min (remakes) and those with early
disconnects are excluded. Each match yields two team-level observations
($N_{\text{obs}} = 4{,}000$).

\textbf{Variables.}

\vspace{2pt}
\begin{tabular}{@{}p{3.6cm} p{10.8cm}@{}}
\toprule
\textbf{Variable} & \textbf{Operationalization} \\
\midrule
\textit{Win} (DV) & Binary: 1 = team won, 0 = lost. \\
\textit{Toxicity Score} (IV1) & Standardized composite of: (a)~\textbf{Feeding Index}
  — $\max_i \frac{\text{deaths}_i}{\text{kills}_i+\text{assists}_i+1}$ across team;
  (b)~\textbf{Early Surrender} — binary, surrender vote before min~20;
  (c)~\textbf{Vision Neglect} — share of players with vision score below 10th
  percentile for role. \\
\textit{Leadership Score} (IV2) & Standardized composite of:
  (a)~\textbf{Coordinating Ping Ratio} — share of coordinating pings
  (onMyWay + command + assistMe + push) in total pings for top-pinging player;
  (b)~\textbf{Objective Coordination} — count of major objectives taken with
  all five members alive and nearby;
  (c)~\textbf{Vision Leadership} — support \& jungler vision score above
  role-median. \\
\textit{Controls} & Rank tier, champion-composition type, team-average KDA,
  team-average CS/min, patch fixed effects. \\
\bottomrule
\end{tabular}

\vspace{4pt}
\textbf{Model.} Binary logistic regression, estimated stepwise:

\vspace{-4pt}
{\small
\[
\textstyle
\log\!\frac{P(\text{Win})}{1-P(\text{Win})}
= \beta_0 + \beta_1\,\text{Tox} + \beta_2\,\text{Lead}
+ \beta_3\,(\text{Tox}\!\times\!\text{Lead})
+ \boldsymbol{\beta_C}\mathbf{X_C}
\]}

\vspace{-2pt}
Robustness: (i)~OLS on Gold Differential at 15~min; (ii)~stratified by rank tier;
(iii)~sub-components tested individually. Data collection via Python
(\texttt{riotwatcher}); analysis in R (\texttt{glm}).

% ============================================================
\section{Expected Contribution}
% ============================================================

This study provides the first operationalization of leadership-style communication
in LoL using publicly accessible ping data — a construct no prior paper has
tested. It complements the toxicity literature by modeling behavioral proxies at
scale with proper skill controls, and the interaction term tests a theoretically
grounded moderation hypothesis. Findings inform esports coaching practice and
Riot's behavioral incentive design.

\vspace{6pt}
\hrule
\vspace{4pt}

% ---------- References ----------
{\small
\printbibliography[heading=none]
}

\end{document}
